\documentclass[14pt,a4paper,oneside]{article}

\usepackage[utf8]{inputenc}   


\usepackage{amsmath,amsthm,amssymb,appendix,bm,mathrsfs,ctex,geometry,tcolorbox,fancyhdr,tikz,esint,appendix}


%% GB/T 7714 宏包 ---- 引文格式
\usepackage[sort&compress]{gbt7714}
\bibliographystyle{gbt7714-numerical}



% \numberwithin{equation}{section}             %%%%%% 公式编号




\geometry{left=1.7cm, right=1.7cm, top=2.54cm, bottom=2.54cm}    %%%%% 页面尺寸
\linespread{1.5}                             %%%%% 行距1.5倍
\pagestyle{fancy}



%%    设置页眉页脚
\fancyhf{}
\renewcommand{\headrulewidth}{0.15mm}
\fancyhead[C]{文献阅读报告}                   %%%%%% 页眉
\fancyfoot[C]{\thepage}                      %%%%%% 页脚 页码
\setlength{\headheight}{15pt}


%%   设置各种链接颜色为黑色
\usepackage[colorlinks,
            linkcolor=black,
            anchorcolor=black,
            citecolor=black,
            urlcolor=blue
            ]{hyperref}



%%   例子编号
\newcounter{example}
\counterwithin{example}{subsection}
\newcommand{\examplename}{Example}
\newcommand{\exampleTitle}[1]{%
\textbf{\refstepcounter{example}\examplename\theexample:#1}\par}


%%  插入图片配置
\usepackage{graphicx,subfigure}
\usepackage[]{caption2} %新增调用的宏包
% \numberwithin{figure}{section}               %%%%%% 图片编号按节

\graphicspath{ {./Figure/} }            %%%%% 设置图片文件所在目录
\DeclareGraphicsExtensions{.pdf,.svg,.jpeg,.png,.jpg}     %%%%% 支持图片格式


%%    支持tikz作图
\usepackage{tikz}
\usetikzlibrary{backgrounds,arrows}


\usepackage{pgfplots}
\pgfplotsset{compat=1.15}
\usepackage{mathrsfs}


\title{\textbf{文献阅读报告 13\\Millijoule few-cycle pulses from staged compression for strong and high field science}\cite{stanfieldMillijouleFewcyclePulses2021}}                           %%%%% 标题
\author{骆禹天}                                           %%%%% 作者
\date{2024/**/**}                                        %%%%%% 日期

%%   定义一些数学格式
\newtheorem{theorem}{定理}[]
\newtheorem{definition}[theorem]{定义}
\newtheorem{lemma}[theorem]{引理}
\newtheorem{corollary}[theorem]{推论}
\newtheorem{example}[theorem]{例}
\newtheorem{proposition}[theorem]{命题}
\renewcommand{\abstractname}{\Large\textbf{摘要}}




%------------------------- 开始文档  -------------------------%



\begin{document}

\pagenumbering{Roman}         %%%%% 目录前页码统计按罗马数字

\maketitle

\setcounter{page}{1}          %%%%% 标题开始统计页数



%------------------------- 目录  ------------------------------%

\tableofcontents              %%%%% 插入目录

\newpage
\setcounter{page}{1}          %%%%% 目录后重新统计页数


\pagenumbering{arabic}        %%%%% 目录后统计页码为阿拉伯数字



%%%%%%%%%%%%%%%%%%%% 开始第一节  %%%%%%%%%%%%%%%%%%%%%


\section{历史}

1834年, Russell 观察到了孤子现象:
> 1. 他观察到孤立波，这是一种长而浅的永久形式的水波，因此他推断出它们的存在；这是他最重要的成果； 
> 2. 孤立波在深度 $h$ 均匀的通道中的传播速度 $c$ 由以下公式给出：$c^2 = g(h +\eta )$，其中, $\eta$ 是波的振幅，$g$ 是重力。


Korteweg-de Vries (KdV) 方程

控制在浅水通道中传播的长一维、小振幅、表面重力波的非线性演化方程
\begin{equation}
   \frac{\partial\eta}{\partial\tau}=\frac32\sqrt{\frac gh}\frac\partial{\partial\xi}\left(\frac12\eta^2+\frac23\alpha\eta+\frac13\sigma\frac{\partial^2\eta}{\partial\xi^2}\right),\quad\sigma=\frac13h^3-Th/(\rho g)
\end{equation}


其中,
- $\eta$ 是波在平衡水平 $h$ 之上的表面高度，
- $a$ 是与液体匀速运动相关的任意小常数，
- $g$ 是引力常数，
- $T$ 是表面张力，
- $\rho$ 是密度
- （术语“长”和“小”是指与通道的深度相比）。

方程（1.1.1）可以通过变换
\begin{equation}
    t=\frac{\mathrm{l}}2\sqrt{g/(h\sigma)}\:\tau,\quad x=-\sigma^{-1/2}\xi,\quad u=\frac{\mathrm{l}}2\eta+\frac{\mathrm{l}}3\alpha 
\end{equation}

转化为无量纲形式
\begin{equation}
    u_t+6uu_x+u_{xxx}=0
\end{equation}

其中, 下标表示偏微分.

KdV 方程具有孤立波解:
\begin{equation}
    u(x,t)=2\kappa^2\sinh^2\left\{\kappa(x-4\kappa^2t-x_0)\right\}\tag{1.1.3}
    \label{eq:}
\end{equation}



%--------------------  插入参考文献  --------------------%

\newpage

\clearpage           %%%  “目录”中显示的页码正确, 需要添加 \clearpage
\phantomsection      %%   “目录”中的链接能正确跳转, 需要添加 \phantomsection
\addcontentsline{toc}{section}{参考文献}   %向目录中添加条目, 以章的名义

\bibliography{ref}

\end{document}